\emph{Setting and tools.} Throughout, $(M,G)$ satisfies (T1) to (T3) of \claimref{sheaf.tameness}, and we use the definitions and clauses (1) to (3) of \claimref{synthesis.wall-rung-count-is-height-of-dying-line-class}: a $k$-rung completion is a smooth symmetric form $h$ on $T_{k} = T^{*}M \oplus \underline{\R}^{2k}$ over $\Mp$, nondegenerate on $\Mp$ including $\Sig$, with $h|_{T^{*}M} = G$ along $\Sig$ and $h = G \perp H_{k}$ outside the collar; $\overline G$ is the induced form on $T^{*}M|_{\Sig}/\shL$, of signature $(p-1, q)$ (\claimref{sheaf.degeneration-sheaf}), and $N(\cdot)$ denotes a maximal negative subbundle, so $N(\overline G)$ has rank $q$. Line bundles with equal $w_{1}$ are isomorphic (Fact \claimref{sheaf.imported-real-line-bundles}: tensor squares of line bundles are trivial), and $w$ is multiplicative, $w(E \oplus E') = w(E)\,w(E')$ (Fact \claimref{synthesis.imported-whitney-product-formula-and-projective-space}(i)). A bundle $E$ is \emph{stably trivial} if $E \oplus \underline{\R}^{r} \cong \underline{\R}^{r + \operatorname{rank} E}$ for some $r$. The rotor double cover $\mathrm{Spin}(n) \to \mathrm{SO}(n)$, $R \mapsto (v \mapsto R v \tilde R)$, is \claimref{std-algebra.spin-double-cover}.

Classical inputs, cited from memory and not in the store; every step that uses them is conditional on them, and deriving them is an open item. (L1) For $\Sig_{m} = S^{2m+1}/\langle J \rangle$ (below): $H^{2j}(\Sig_{m}; \mathbb{Z}) \cong \mathbb{Z}_{4}$ generated by $y^{j}$ for $1 \le j \le m$, with $y$ the Euler class of the quarter-turn plane bundle $\eta$ (the Gysin sequence of the circle bundle $\Sig_{m} \to \C P^{m}$); $H^{*}(\Sig_{m}; \mathbb{Z}_{2}) = \Lambda[x] \otimes \mathbb{Z}_{2}[u]/(u^{m+1})$ with $u = y \bmod 2$; and $T\Sig_{m} \oplus \underline{\R} \cong (m+1)\,\eta$, so $w(\Sig_{m}) = (1 + u)^{m+1}$ (the tangent bundle of $S^{2m+1}$ plus its normal line is $\R^{2m+2}$, on which the deck transformation acts by $J$). (L2) $H^{4}(B\mathrm{Spin}; \mathbb{Z}) \cong \mathbb{Z}$ is generated by the spin class $q_{1}$, natural and stable (unchanged by adding a trivial summand with the induced spin structure), with $2q_{1} = p_{1}$ (E.~Thomas); $p_{1}(P) = e(P)^{2}$ for an oriented plane bundle $P$ (Milnor and Stasheff, Cor.~15.8); and a bundle is orientable iff $w_{1} = 0$, spin iff moreover $w_{2} = 0$. (L3) For the cross-check only: $K(\Sig_{m}) = \mathbb{Z}[\eta]/(\eta^{4} - 1, (\eta - 1)^{m+1})$ (Kambe 1966; Atiyah; from the Thom isomorphism of the free $\mathbb{Z}_{4}$ action on $S^{2m+1}$). The same computation at $m = 3$ is clause (1) at $n = 3$ of a sibling K-theoretic result (named in the extraction note).

\emph{(1).} Let $h$ be a $k$-rung completion with $k \le 2$. If $k = 1$, then $x = 0$ (clause (3) of the parent), so $\shL$ is trivial and so is $\shL \oplus \shL$. If $k = 2$, clause (1) of the parent gives a line bundle $N$ over $\Sig$ with
\[
N(\overline G) \oplus \shL \oplus N \;\cong\; N(\overline G) \oplus \underline{\R}^{2}, \qquad w(N) = (1 + x)^{-1}.
\]
Since $N$ is a line, $w(N) = 1 + w_{1}(N)$, so $w_{1}(N) = x$ (and $x^{2} = 0$), and $N \cong \shL$. Choose a bundle $K$ with $N(\overline G) \oplus K \cong \underline{\R}^{r}$ ($\Sig$ is compact); adding $K$ to both sides gives $\shL \oplus \shL \oplus \underline{\R}^{r} \cong \underline{\R}^{r+2}$. The negative part $N(\overline G)$ of the physical metric, which the indefinite completion on $L(4,1)$ uses, cancels in this step: it can change the isomorphism class, never the stable class. If $q = 0$, then $N(\overline G) = 0$ and the displayed isomorphism with $k$ rungs (clause (1) of the parent for general $k$, $N$ of rank $k - 1$) reads $\shL \oplus N \cong \underline{\R}^{k}$, so $\shL$ embeds in $\underline{\R}^{k}$ and $e(\shL) \le k(\Sig)$; with clause (3) of the parent, $k(\Sig) = e(\shL)$.

\emph{(2).} An oriented plane bundle $P$ is a bundle with structure group $\mathrm{SO}(2)$, and a spin structure is a lift through $\mathrm{Spin}(2) \to \mathrm{SO}(2)$, which in GA is $R = \exp(e_{1}e_{2}\,\varphi) \mapsto (v \mapsto R v \tilde R)$, the rotation by $2\varphi$: the squaring map of the circle. By the Whitney formula $w(\shL \oplus \shL) = (1 + x)^{2} = 1 + x^{2}$, so if $x^{2} = 0$ then $w_{1}(\shL \oplus \shL) = 0$ and $w_{2}(\shL \oplus \shL) = 0$, and by (L2) $\shL \oplus \shL$ is orientable and spin. On the classifying space, $H^{*}(B\mathrm{Spin}(2); \mathbb{Z}) = \mathbb{Z}[a]$ and the Euler class pulls back to $\pm 2a$ (squaring has degree two), so $p_{1} = e^{2} = 4a^{2}$ by (L2), and since $H^{4}(B\mathrm{Spin}(2); \mathbb{Z})$ is torsion-free, $q_{1} = 2a^{2}$. By naturality, for a spin plane bundle whose $\mathrm{Spin}(2)$ bundle, read as an oriented plane bundle through the isomorphism $\exp(e_{1}e_{2}\varphi) \mapsto$ rotation by $\varphi$, has Euler class $a$,
\[
q_{1} = 2a^{2}, \qquad p_{1} = 4a^{2}, \qquad e = 2a .
\]
Now suppose $\shL \oplus \shL \oplus \underline{\R}^{r} \cong \underline{\R}^{r+2}$, $r \geq 1$; composing with a reflection of one trivial summand we may take the isomorphism orientation preserving. Spin structures on $\shL \oplus \shL$ and on $\shL \oplus \shL \oplus \underline{\R}^{r}$ correspond under stabilisation, and the isomorphism carries them bijectively to the spin structures of $\underline{\R}^{r+2}$, among which is the product structure, with $q_{1} = 0$ (it is pulled back from a point). By stability (L2), some spin structure of $\shL \oplus \shL$ has $q_{1} = 0$. Contrapositively, if every spin structure of $\shL \oplus \shL$ has $q_{1} \neq 0$, then $\shL \oplus \shL$ is not stably trivial, and by (1) $k(\Sig) \geq 3$ on every side.

\emph{(3), the lens walls.} In $\Cl(2m+2, 0)$ let $B = e_{1}e_{2} + \dots + e_{2m+1}e_{2m+2}$ and $J(v) := v \cdot B = \tfrac12 (vB - Bv)$, the quarter turn $e_{2j-1} \mapsto e_{2j} \mapsto -e_{2j-1}$ in each coordinate plane. $J$ is orthogonal with $J^{2} = -1$, so $J^{2}x = -x \neq x$ for $x \neq 0$ and $\langle J \rangle \cong \mathbb{Z}_{4}$ acts freely on the unit sphere; the quotient is $\Sig_{m} = L^{2m+1}(4)$. Sections of $\shL$ are functions with $f(Jx) = -f(x)$; $\eta$ is the flat plane bundle on which $J$ acts by the quarter turn $J_{2}(v) = v \cdot e_{1}e_{2}$ of $\R^{2}$. By (L1), $x = w_{1}(\shL) \neq 0$ (the holonomy is $-1$) and $x^{2} = 0$, so $\operatorname{ht}(x) = 2$.

The doubled line $\shL \oplus \shL = \shL \otimes \R^{2}$ is the flat plane bundle on which $J$ acts by $-1$, the half turn, which is the image of the rotor $e_{1}e_{2}$ under the double cover: $e_{1}e_{2}\, v\, (e_{1}e_{2})^{\sim} = -v$ on the plane, $(e_{1}e_{2})^{2} = -1$. The lifts of this holonomy to $\mathrm{Spin}(2)$ are $\pm e_{1}e_{2}$, both of order four, and they give the two spin structures of $\shL \oplus \shL$ (spin structures form a torsor over $H^{1}(\Sig_{m}; \mathbb{Z}_{2}) = \mathbb{Z}_{2}$, and the two flat lifts differ by the character $J \mapsto -1$, which is $x \neq 0$). Under $\exp(e_{1}e_{2}\varphi) \mapsto$ rotation by $\varphi$, the rotor $e_{1}e_{2} = \exp(e_{1}e_{2}\,\pi/2)$ goes to the quarter turn $J_{2}$, and $-e_{1}e_{2} = (e_{1}e_{2})^{3}$ to $J_{2}^{3} = -J_{2}$; so the associated plane bundles are $\eta$ and its conjugate, $a = \pm y$ (equivalently $y$ or $3y$), and by (2),
\[
e(\shL \oplus \shL) = 2y \neq 0, \qquad p_{1}(\shL \oplus \shL) = 4y^{2} = 0, \qquad q_{1}(\shL \oplus \shL) = 2y^{2}
\]
for both spin structures. By (L1), $2y^{2} \neq 0$ in $H^{4}(\Sig_{m}; \mathbb{Z}) = \mathbb{Z}_{4}$ when $m \geq 2$ (for $m = 1$, $H^{4} = 0$). By (2), $k(\Sig) \geq 3$ for every tame wall with cross-section $\Sig_{m}$, $m \geq 2$, and dying line $\shL$, on every side; note that $p_{1}$ does not see this.

\emph{Existence.} On $M = \Sig_{m} \times \R$, identify $T^{*}M|_{\Sig_{m}} = T^{*}\Sig_{m} \oplus \R\,dt$ with the bundle $(S^{2m+1} \times \R^{2m+2})/\langle J \rangle$, the deck transformation acting by $J$ on both factors (round metric, normal line, (L1)). Let $c$ be the linear map negating $e_{2}, e_{4}, \dots, e_{2m+2}$; then $cJ = -Jc$, so $x \mapsto c(x)$ satisfies $c(Jx) = -J\,c(x)$, is nowhere zero, and is a section of $\shL \otimes T^{*}M|_{\Sig_{m}}$ without zeros: it embeds $\shL$ as the line spanned by $c(x)$. With $\pi_{x}$ the Euclidean projector onto that line, $\pi_{Jx} = J \pi_{x} J^{T}$, so the cometrics
\[
G_{t} := t\,\pi - (1 - \pi) \quad (\text{indefinite}), \qquad G_{t} := 1 - (1 - t)\,\pi \quad (\text{definite})
\]
descend to $M$. Each has eigenvalue $t$ on $\shL$ and $\mp 1$, respectively $1$, on its complement: $\det G_{t} = -t$, respectively $t$, with a simple zero, corank one exactly at $t = 0$ with radical $\shL$, so (T1) to (T3) hold as for the projective walls of the parent claim (Jacobi's formula); the physical side $t > 0$ has signature $(1, 2m+1)$, respectively $(2m+2, 0)$. On $\Sig_{m} \times T^{j}$ take $G_{t} \perp A$ with $A$ a constant definite form on $T^{*}T^{j}$; the dying line and its classes are pulled back along the projection, which has a section, so $x \neq 0 = x^{2}$ and $q_{1}(\shL \oplus \shL) \neq 0$ persist, and $k \geq 3$ in every dimension $2m+1+j \geq 5$.

\emph{The upper bound on $L^{7}(4)$.} By clause (2) of the parent, $k(\Sig) \le e(\shL)$ on every side. On $\R^{8} = \Cl^{0}(3,0)^{2}$, pairs $(R_{1}, R_{2})$ of even elements, right multiplication by $e_{2}e_{3}$ is an orthogonal complex structure $J'$; any two orthogonal complex structures on $\R^{8}$ are conjugate by an orthogonal map (each has an orthonormal basis $u_{1}, J u_{1}, \dots, u_{4}, J u_{4}$), so $S^{7}/\langle J' \rangle$ is isometric to $\Sig_{3}$ with $\shL$ corresponding to $\shL$. Since $e_{2}e_{3}\, e_{2}\, (e_{2}e_{3})^{\sim} = -e_{2}$, the vector $R e_{2} \tilde R$ is odd under $R \mapsto R e_{2}e_{3}$, and $|R e_{2}\tilde R|^{2} = (R\tilde R)^{2}$. So $F = (R_{1}e_{2}\tilde R_{1}, R_{2} e_{2} \tilde R_{2}) \in \R^{6}$ is odd and $|F|^{2} = (R_{1}\tilde R_{1})^{2} + (R_{2}\tilde R_{2})^{2} > 0$ on the sphere: $\shL$ embeds in $\underline{\R}^{6}$, $e(\shL) \le 6$, and $3 \le k(\Sig_{3}) \le 6$ on every side.

\emph{(4), dimension five.} Take $m = 2$ and the indefinite cometric of the existence paragraph: a tame wall on $L^{5}(4) \times \R$ with physical side of signature $(1,5)$ and $\operatorname{ht}(x) = 2$. By (3), $k(\Sig) \geq 3$. The mod-2 data cannot see the difference ($x \neq 0 = x^{2}$, as on the torus and lens walls of dimension three where an indefinite side needs two), nor can the Euler class $\beta x = 2y$ through its square ($4y^{2} = 0$, so $p_{1} = 0$); a necessary condition for two rungs is that $\shL \oplus \shL$ be stably trivial, and on $L^{5}(4)$ it fails, detected by $q_{1} = 2y^{2}$ (sufficiency is not proved). Note that $q_{1} = 2a^{2}$ for any $a$ with $2a = \beta x$ (remark below), so $q_{1}$ is a function of $\beta x$ through the integral cup product; what vanishes is $(\beta x)^{2} = p_{1}$. This result therefore does not show that $k(\Sig)$ is undetermined by $W(\Sig)$, $\beta x$ and the signature. For comparison, on $L(4,1)$ ($m = 1$) $H^{4} = 0$, so the obstruction vanishes, consistent with the two-rung completion on indefinite sides of \claimref{synthesis.lens-space-wall-rung-count-sees-the-integral-twist}.

\emph{Checks.} The check verifies in GA and exact arithmetic: $J^{2} = -1$, $J$ orthogonal, $cJ = -Jc$, the equivariance $\pi_{Jx} = J\pi_{x}J^{T}$, $\det G_{t}$, the signatures on both sides and the radical at $t = 0$ for $m = 2, 3$; the spin lift ($e_{1}e_{2}$ squares to $-1$, covers the half turn, corresponds to the quarter turn); the class arithmetic of (2) and (3) in $\mathbb{Z}[a]$ and $\mathbb{Z}[y]/(4y, y^{m+1})$, including $q_{1} = 0$ for $m = 1$ and $w = 1$ exactly for $m = 1, 3$; and the quaternionic odd map.

\emph{Remark (cross-check in K-theory; (L3), from memory).} The complexification of $\shL \oplus \shL$ minus the trivial bundle is $2(\eta^{2} - 1)$ in $K(\Sig_{m})$ (the flat line of holonomy $-1$, complexified, is $\eta^{2}$). With $s = \eta - 1$ this is $4s + 2s^{2}$, and the check shows it lies outside the ideal $((1+s)^{4} - 1, s^{m+1})$ for $m = 3$ (it reduces to $2s^{3} \neq 0$) and inside it for $m = 1, 2$. So on $L^{7}(4)$ complex K-theory confirms independently that $\shL \oplus \shL$ is not stably trivial, while on $L^{5}(4)$ it is blind and the obstruction is the integral class $q_{1}$ alone (argued: complexification doubles the class on the four-dimensional stage, $2 \cdot 2y^{2} = 0$).

\emph{Remark (the Adams lead; argued).} On $\R P^{n}$ the only nontrivial line is the tautological one, with $x^{n} \neq 0$, so $\operatorname{ht}(x) = e(\shL) = n+1$ and the parent claim already gives $k = n + 1$: real projective walls cannot separate the bounds, and the order $2^{\varphi(n)}$ of $\shL - 1$ in real K-theory gives only $k \geq 1 + \varphi(n)$, weaker than $n + 1$.

\emph{Remark (dimension four and the converse; argued, classical inputs from memory).} The map $B\mathrm{Spin} \to K(\mathbb{Z}, 4)$ given by $q_{1}$ is $8$-connected (Bott periodicity: $\pi_{i}(\mathrm{SO}) = 0$ for $i = 4, 5, 6$), so over a complex of dimension at most seven a spin bundle is stably trivial if and only if $q_{1} = 0$. For a general wall with $x^{2} = 0$, the Bockstein sequence gives $\beta x = 2a$ and $\shL \oplus \shL$ is the square of the plane bundle with Euler class $a$, so $q_{1} = 2a^{2}$, independent of the choice of $a$. On a closed four-manifold $a$ is torsion and $H^{4}(\Sig; \mathbb{Z})$ is $\mathbb{Z}$ or $\mathbb{Z}_{2}$, so $2a^{2} = 0$: in dimension four the stable condition of (1) is implied by $x^{2} = 0$, and whether $k(\Sig) = \operatorname{ht}(x)$ on indefinite four-dimensional walls is not decided here. Dimension five is the first in which the stable class of $\shL \oplus \shL$ can exceed the mod-2 data, and $L^{5}(4)$ realises it.
